\subsection{Entraînement}\label{sub:Entraînement} % (fold)

%usage pragmatique du tool en cli -> révision de la stratégie d'entrainement et de la fonction du prédicteur
%-> focalisation sur a prédiction de l'embedding: 
%   d'un token aléatoire
%du dernier tokendu dernier token, mais avec le dernier tier des tokens d'entrée masqués, pour forcer l'anticipation
%-> tentative de restreindre la dimension de l'embedding (256 -> 128 -> 64) pour faire converger les représentation


\subsubsection{Infrastructure et configuration}

L'entraînement est piloté via un fichier de configuration YAML, garantissant la
reproductibilité complète de chaque run. Trois runs successifs sont archivés sur le dépôt
HuggingFace \texttt{rubenanko/binary-jepa}, chacun stockant ses checkpoints par époque
(\texttt{ep\{n\}.pt}), le checkpoint de référence (\texttt{latest.pt}), le journal de
perte (\texttt{log.csv}) et sa configuration (\texttt{train.yaml}). Le corpus d'entraînement
est \texttt{raw\_data500.json}, un sous-ensemble des séquences extraites par le pipeline
de la Section~\ref{sub:Extraction et Représentation des Données}.

\subsubsection{Historique des runs et analyse de convergence}
\label{sub:runs}

Le Tableau~\ref{tab:runs} résume les configurations et résultats des trois runs conduits.

\begin{table}[h]
\centering
\caption{Récapitulatif des runs d'entraînement archivés sur \texttt{rubenanko/binary-jepa}.}
\label{tab:runs}
\begin{tabular}{lccccc}
\toprule
Run & Date & Stratégie de masquage & LR & Batch & Perte finale \\
\midrule
\texttt{10-03-2026} & 10 mars 2026 & Tiers consécutifs & $5 \times 10^{-5}$ & 20 & 512.88 \\
\texttt{last-token} & 14 mars 2026 & Dernier token & $10^{-4}$ & 10 & 5.81 \\
\texttt{random-token} & 14 mars 2026 & Token aléatoire & $10^{-4}$ & 10 & 0.030 \\
\bottomrule
\end{tabular}
\end{table}

\paragraph{Run 1 — \texttt{10-03-2026} : divergence catastrophique.}
Le premier run ($\eta = 5 \times 10^{-5}$, $B = 20$, stratégie de masquage par tiers
consécutifs) présente une divergence dès l'époque 1 : la perte Huber atteint 14.82 à la
première époque, puis croît de façon monotone pour atteindre 512.88 à l'époque 10.
Une telle trajectoire exclut tout problème de taux d'apprentissage trop élevé — la perte
n'est jamais décroissante. Ce comportement est cohérent avec un effondrement de
représentation (\textit{representation collapse}) survenant avant même la première mise à
jour EMA stabilisante, vraisemblablement amplifié par la taille de batch élevée ($B=20$)
réduisant le bruit stochastique protecteur dans les premières itérations. Ce run est
abandonné.

\paragraph{Run 2 — \texttt{last-token} : instabilité tardive.}
Le second run ramène la taille de batch à $B = 10$ et adopte une stratégie de masquage
dit \textit{dernier token} : seul le token terminal de la séquence est masqué, et c'est
sa représentation latente que le prédicteur $g_\varphi$ doit reconstruire. La perte
s'établit à des valeurs raisonnables lors des premières époques
($\mathcal{L}_1 = 0.672$, $\mathcal{L}_7 = 0.997$), avant de diverger à partir de
l'époque 8 ($\mathcal{L}_8 = 2.256$, $\mathcal{L}_{10} = 5.807$). Cette instabilité
tardive suggère que la tâche du dernier token est insuffisamment contraignante :
l'encodeur apprend à maximiser la prédictibilité d'une unique position terminale,
au risque de surspécialiser les représentations et de rendre l'encodeur cible $f_\xi$
trop sensible aux derniers tokens d'API ou de JumpKind, qui sont pourtant les plus
variables selon le contexte d'obfuscation. Ce run est écarté au profit du run suivant.

\paragraph{Run 3 — \texttt{random-token} : convergence stable.}
Le troisième run conserve la configuration $B = 10$, $\eta = 10^{-4}$, et adopte un
masquage aléatoire des positions cibles au sein de la séquence réelle (hors padding).
La courbe de perte converge régulièrement :

\begin{equation}
\mathcal{L}_1 = 0.0776 \;\to\; \mathcal{L}_3 = 0.0375 \;\to\; \mathcal{L}_7 = 0.0303
\;\to\; \mathcal{L}_{10} = 0.0300,
\end{equation}

indiquant un plateau atteint dès l'époque 7 ($\Delta\mathcal{L}_{7 \to 10} < 10^{-3}$).
Ce plateau constitue un signal de convergence sain dans le cadre JEPA : il ne signifie
pas que le modèle a cessé d'apprendre, mais que l'encodeur cible $f_\xi$ et l'encodeur
de contexte $f_\theta$ ont atteint un équilibre dynamique stable sous la contrainte EMA.
Le checkpoint \texttt{random-token/latest.pt} — correspondant à l'époque 10 — est promu
comme artefact de référence et copié à la racine du dépôt (\texttt{latest.pt}), constituant
le modèle utilisé dans les expérimentations de la Section~\ref{sec:experiments}.

\subsubsection{Discussion : impact de la stratégie de masquage}

Les trois runs révèlent que la stratégie de masquage est le facteur déterminant de la
convergence, devant le taux d'apprentissage ou la taille de batch. Le masquage par tiers
consécutifs (run 1) impose une tâche de prédiction de continuation à longue portée que
l'encodeur convolutif — dont le champ réceptif effectif est limité à 13 tokens — ne
peut structurellement résoudre dans les premières itérations, conduisant à l'explosion du
gradient. Le masquage du dernier token seul (run 2) réduit trop la difficulté de la
tâche et produit des représentations surspécialisées. Le masquage aléatoire (run 3)
offre un compromis adapté : la tâche est localement résoluble par le champ réceptif
convolutif, mais suffisamment diverse pour empêcher la mémorisation et forcer la
généralisation sémantique.

Il convient de noter que la stratégie de masquage par tiers consécutifs décrite en
Section~\ref{sub:masking} constitue l'objectif d'entraînement cible de l'architecture
I-JEPA 1D, mais n'a pas encore produit de run convergent dans le cadre de cette preuve
de concept. Son application requiert vraisemblablement un encodeur à champ réceptif
global (Transformer, Mamba) ou un curriculum d'entraînement progressif augmentant
graduellement la longueur du bloc masqué.

% subsection Entraînement (end)
